%mscS --> seminar
%msc --> Thesis

% برای گزارش سمینار mscS  و برای پایان‌نامه ,msc را انتخاب کنید
\documentclass[oneside,openany,msc]{SBU-Thesis}

% برای وارد کردن بسته‌هایی که نیاز است قبل از bidi وارد شوند، آنها را در فایل confix.tex قرار دهید. سایر بسته‌های مورد نیاز را در اینجا می‌توانید وارد کنید:

\begin{document}
	%عنوان پایان‌نامه
	\title{راه‌کار‌ی مبتنی بر یادگیری ماشین برای پیش‌بینی مهارت‌های افراد در آینده در وبسایت‌های پرسش‌‌وپاسخ اجتماعی }	
	%رشته
	\subject{مهندسی کامپیوتر}
	
	%گرایش
	\field{نرم‌افزار}
	
	%نام
	\name{سینا}

	%نام خانوادگی
	\surname{محمدی}
	
	%استاد راهنما
	\firstsupervisor{دکتر محمود نشاطی}
	
	%استاد راهنمای دوم (در صورت وجود، در غیر این صورت این خط را حذف کنید)
%	\secondsupervisor{دکتر محمود نشاطی}

	%استاد مشاور (در صورت وجود، در غیر این صورت این خط را حذف کنید)
%	\firstadvisor{دکتر محمود نشاطی}

%تاریخ انجام پایان‌نامه
	\thesisdate{تابستان ۱۴۰۴}
	
	%نام دانشکده
	\faculty{دانشکده مهندسی و علوم کامپیوتر}
	
	%نام دانشگاه
	\university{دانشگاه شهید بهشتی}
	
	%اسامی برای صفحات داوران: در هر مورد، در صورت عدم وجود، کل خط را حذف کنید
	\davaranSupervisor{نام و نام‌خانوادگی}
	\davaranSecondSupervisor{نام و نام‌خانوادگی}
	\davaranAdvisor{نام و نام‌خانوادگی}
	\davaranInternal{نام و نام‌خانوادگی}
	\davaranExternal{نام و نام‌خانوادگی}
	\davaranAssignee{نام و نام‌خانوادگی}
%	\davaranDate{۱۳۹۷/۱۱/۷} % چنانچه این خط را حذف کنید با فضای خالی جایگزین می‌شود

	%مقادیر انگلیسی برای صفحه آخر
	\latinuniversity{Shahid Beheshti University}
	\latinfaculty{Faculty of Computer Science \& Engineering}
	\latintitle{A Machine Learning-based Approach for Predicting Future Skills of Individuals on Community Question Answering websites}
	\latinname{sina}
	\latinsurname{mohammadi}
	\latinthesisdate{2025}
	\firstlatinsupervisor{Dr. Mahmood Neshati}
	%	\secondlatinsupervisor{Dr. Mahmood Neshati}
	%	\firstlatinadvisor{Dr. Mahmood Neshati}
	% چکیده به فارسی
	\fa-abstract{چکیده در اینجا قرار میگیرد}
	%کلمات کلیدی:
	\keywords{}
	%%%%%%%%%%%%%%%%%%%%%%%%%%%%%%%%%
	
	
\firstPage %ساخت صفحه اول پایان‌نامه 
\davaranPage % ساخت صفحه امضا داوران

\rightsPage % ساخت صفحه حقوق پایان‌نامه
\copyRightPage % ساخت صفحه تعهدنامه

%%%%%%%%%%%%%%%%%%%%%
\abstractPage % ساخت صفحه چکیده


\tableofcontents % فهرست مطالب
\listoffigures \newpage % فهرست تصاویر
\listoftables \newpage % فهرست جداول

		

%%%%%%%%%%%%%%%%%%%%%%%%%%%

\include{body}

%%%%%%%%%%%%%%%%%%%%%%%%%%%

% مراجع
\newpage
%\onehalfspacing
\bibliographystyle{ieeetr-fa}%{chicago-fa}%{plainnat-fa}%
\bibliography{thesis}


%%%%%%%%%% نمونه‌ی واژه‌نلمه تک ستونه ٪٪٪٪٪٪٪٪٪٪
\baselineskip=.6cm

\chapter*{واژه‌نامه  انگلیسی به  فارسی}\markboth{واژه‌نامه  انگلیسی به  فارسی}{واژه‌نامه  انگلیسی به  فارسی}
%\thispagestyle{fancy}
\addcontentsline{toc}{chapter}{واژه‌نامه  انگلیسی به  فارسی}
\noindent
\englishgloss{Word}{کلمه}
\englishgloss{Word 2}{کلمه ۲}
\englishgloss{Word 3}{کلمه ۳}


\chapter*{واژه‌نامه فارسی به انگلیسی}\markboth{واژه‌نامه فارسی به انگلیسی}{واژه‌نامه فارسی به انگلیسی}
\addcontentsline{toc}{chapter}{واژه‌نامه فارسی به انگلیسی}
%\thispagestyle{fancy}
\noindent
\persiangloss{کلمه ۱}{word 1}
\persiangloss{کلمه ۲}{word 2}
\persiangloss{کلمه ۳}{word 3}

\baselineskip=1cm
%%%%%%%%%%%%%%%%%%%%%%%%%


% چکیده به انگلیسی
\en-abstract
{
	This is Abstract in English.
}

% کلمات کلیدی انگلیسی 
\latinkeywords
{
	Word1, Word2, Word3
}

\latinAbstractPage % ساخت صفحه چکیده به انگلیسی
\latinFirstPage % ساخت صفحه آخر

	
\end{document}